\documentclass{article}

% This is a comment line in latex

\usepackage[T1]{fontenc}
\usepackage[utf8]{inputenc}
\usepackage[french]{babel}

%% This package is necessary to use \includegraphics.
\usepackage{graphicx}

%% This package is necessary to define hyperlinks.
\usepackage{hyperref}

%% These packages are necessary to include code.
\usepackage{listings}
\usepackage{minted} % colored

%% This package is needed to enchance mathematical formulas.
\usepackage{amsmath}


% Latex allows you to define your own "commands",
% better known as "macros" in the Latex world.
% The following line is an example of such definition.
\newcommand{\latex}{\LaTeX}

\selectlanguage{french}

\title{%
  Rapport intermédiaire pour le ``Projet de programmation''\\
  Pour la date du 15 Mars, 2025%
}

\author{Cotrez Jules et Shay Ronen}

\begin{document}
\maketitle

\section{Introduction}

\subsection{But du projet}
Le but de notre projet est d'implémenter dans un langage fonctionnel (ici Ocaml) un jeu de plateformes en deux dimensions, dans lequel l'objectif
est de réaliser des niveaux consécutifs en combattant des ennemis ainsi qu'en s'aidant d'items dispersés un peu partout dans le jeu.

\subsection{Métrique}
Nous considérerons notre projet comme réussi si nous parvenons à réaliser un jeu complet avec des déplacements (au minimum les déplacements classiques ,
un grappin qui permet un mouvement de balancier, et un jetpack), plusieurs niveaux, des ennemis et des items divers. Le reste (jeu en réseau, IA pour les 
ennemis, animation pour les attaques, ...) sera considéré comme du bonus.

\section{Implémentation}

\subsection{réalisation du projet}
Nous réalisons le projet en nous attribuant chacun une des tâches du planning (pas de branche séparée pour le moment). Nous utilisons le système de
compilation Dune afin de gérer facilement l'organisation des fichiers et les dépendances. Pour l'interface graphique, nous nous servons de la librairie Raylib.

\subsection{partie du logiciel codée}
Pour le moment nous avons codé les déplacements du joueur (grappin inclus), l'affichage d'un menu ainsi que du premier niveau (avec plateformes et un ennemi),
le parsing de fichiers Json afin de facilement modifier la description des niveaux (nombre et position des plateformes et des ennemis). De plus nous avons
veillé à implémenter la plus grande partie possible de manière fonctionnelle en évitant les effets de bord (c'est là toute la difficultée et l'intérêt de notre
projet !). Les uniques effets de bords proviennent de l'utilisation de librairie Raylib pour l'affichage de l'interface graphique.

\subsection{structure en modules/packages}
Comme dit précédemment, nous avons utilisé le gestionnaire Dune pour instancier la structure de notre projet. Actuellement, nous avons deux classes principales
séparées dans deux fichiers ml (+ deux fichiers mli associés pour l'interface) : jeu et joueur. La classe joueur contient tout ce qui est nécéssaire à l'instaciation
d'un nouveau joueur (joueur principal comme ennemi, nous ferons la distinctions dans l'implémentation du jeu). Elle possède donc tous les attributs utiles tels
que les dimensions, la direction, le vecteur de position, ... La classe jeu contient l'initialisation et la boucle principale du jeu. Elle possède également
des structures facilitant la gestion des différentes entitées (joueurs, ennemis, plateformes).

\end{document}
